Autonomous Surface Vehicles (ASVs) deployed for long-duration environmental monitoring face a fundamental trade-off between information gain and battery constraint. Comprehensive coverage of coastal ecosystems, inland waterways, or offshore regions requires high-fidelity sensors including Radar (50W power), Lidar (30W), optical cameras (5W), and water quality sensors (3W) to capture diverse phenomena from vessel detection to bathymetric mapping to chemical signatures. However, limited onboard battery capacity (typically 10--20 kWh for mid-range ASVs) forces missions to complete within 4--8 hours of continuous operation. Moreover, environmental factors including water currents, wave dynamics, and wind impose variable hydrodynamic drag that scales non-linearly with thrust. An ASV navigating against a 1 m/s current in open ocean may consume 3--5 times more propulsive energy than an equivalent mission in a calm lake, yet no adaptive strategy currently exploits these water-type differences.

\subsection{Problem Motivation}

Long-duration ASV missions encounter three critical challenges in energy-aware multi-sensor exploration. First, sensor-motion coupling creates complex optimization problems. Running all sensors simultaneously maximizes information density but depletes batteries within 2--3 hours, while running minimal sensors extends mission duration at the cost of incomplete coverage and missed events (e.g., hazard detection, species identification). Standard approaches treat sensor scheduling and motion planning independently, missing synergies such as reducing propulsion energy when sensors are inactive or strategically positioning the vehicle to maximize sensor footprint overlap.

Second, water-body type variability demands environment-aware strategies that existing methods ignore. Navigation in a river with 1.5 m/s predictable currents offers opportunities for energy-efficient drift-assisted exploration where the ASV can turn off thrusters periodically and exploit current flow. Coastal environments with 0.5--1.5 m waves demand conservative control to avoid collision and maintain sensor stability, increasing energy expenditure. Open ocean conditions with 2--4 m waves and uncertain current patterns require the most conservative energy budgeting. A one-size-fits-all policy ignores these distinctions, resulting in suboptimal efficiency.

Third, safe return guarantees are difficult to enforce under uncertain battery degradation. Marine batteries degrade due to saltwater exposure, temperature fluctuations, and cycling stress at rates not captured by simple linear models. An ASV that consumes 85\% of its battery reaching a frontier may find itself unable to return safely if degradation effects reduce remaining capacity below predicted reserves. Current rule-based approaches apply conservative fixed thresholds (e.g., return when 50\% battery remains), but these leave significant capacity unexploited in favorable conditions.

Existing solutions prove insufficient for these challenges. Frontier-based exploration methods (RRT*, greedy nearest frontier) maximize coverage but ignore energy consumption entirely. Rule-based energy-aware planning applies fixed duty-cycle schedules (e.g., sensors on for 2 minutes, off for 1 minute) regardless of information value in the current location. Deep reinforcement learning approaches adapted from autonomous driving (DDPG-ASV, Energy-CPP) optimize either coverage or energy efficiency, but not jointly; they treat water-body type as a static nuisance rather than an opportunity for adaptation.

\subsection{Gap Analysis}

We identify a critical gap where no existing method jointly optimizes autonomous surface vehicle motion control and multi-sensor duty-cycling with explicit adaptation to water-body types and guaranteed safe return under battery uncertainty.

Prior work segregates concerns across multiple research threads. Coverage path planning methods such as frontier-based exploration and lawn-mower patterns focus purely on area coverage without energy awareness. Energy-aware path planning including DDPG-ASV and Energy-CPP from agricultural robotics adapts motion to battery state but ignores sensor duty-cycling and assume homogeneous environments. Multi-sensor scheduling in robotics typically addresses static platforms or aerial systems where energy constraints differ fundamentally from marine environments. Reinforcement learning approaches in autonomous systems have not yet addressed the maritime triple-objective of coverage, energy, and multi-modal sensing.

For autonomous surface vehicles specifically, this separation proves suboptimal. Consider a scenario where a policy reaches the 80\% coverage milestone with 25\% battery remaining. A coverage-only baseline might attempt aggressive exploration to reach 90\% coverage, risking battery depletion before returning to base. An energy-only baseline would immediately return to base, leaving 10\% of the mission area unmapped. A jointly-optimized policy would reason about information value of remaining frontiers relative to the energy cost of reaching them and still maintaining safe return—a fundamentally different decision-making process.

Water-body type conditioning has never been incorporated into ASV planning despite its clear importance. A river policy would learn to detect predictable current patterns (via velocity measurements and position history) and exploit them for drift-assisted energy savings. A lake policy would learn that stable conditions permit more aggressive coverage strategies. A coastal policy would learn to balance wave avoidance against information gain. These adaptations cannot emerge from a single policy trained on mixed environments.

\subsection{Proposed Solution}

We propose a deep reinforcement learning system using Soft Actor-Critic (SAC) that jointly optimizes ASV motion control and multi-sensor duty-cycling with explicit water-body type adaptation and energy-constrained exploration. The approach operates across four key dimensions.

First, the system represents mission state as an 87-dimensional vector encoding vehicle dynamics (position, velocity, acceleration to base), battery state (SoC, remaining energy, safe return threshold), environmental conditions (current velocity, wave height, wind, water properties), inferred water-body type (lake, river, coastal, open ocean) via one-hot encoding, coverage progress (frontier locations, information density), sensor states (active sensors, duty cycles), and historical context (recent consumption rates, frontier revisits). This representation enables the policy to reason about causality between state and action such as when currents favor energy-saving drift or when sensor activation will provide high information gain relative to propulsion cost.

Second, a hierarchical action space decomposes exploration decisions into interpretable modes. A discrete Stage 1 selects one of four motion strategies: exploration mode (aggressive coverage, full sensors), energy-saving mode (minimal sensors, current-assisted drift), return-to-base mode (direct path home, minimal sensing), or information-gathering mode (stationary sensing, full sensor suite). Stage 2 provides continuous parameters: thrust percentage, steering angle, speed target, and probabilistic duty cycles for four sensors. This decomposition maintains expressiveness while ensuring human-interpretable decisions and enabling real-time deployment.

Third, a composite reward function balances competing objectives. Coverage reward encourages exploration of new territory. Information reward values sensor measurements weighted by their quality in current conditions (cameras degrade in fog, Lidar degrades in rain). Energy efficiency reward penalizes propulsion consumption. Safety reward enforces a dynamic constraint ensuring sufficient battery reserve for safe return based on real-time distance to base and environmental resistance. These components are weighted to reflect the mission objective: a mapping mission prioritizes coverage and information, while a monitoring mission prioritizes energy efficiency and safe return.

Fourth, the policy conditions explicitly on water-body type to enable environment-specific strategies. During training, missions are conducted across lake (minimal currents, stable conditions), river (1--2 m/s predictable flow), coastal (waves 0.5--1.5 m, tidal currents), and open ocean (waves 2--4 m, variable currents) environments. The policy learns distinct control strategies per water-body type, exploiting current-assisted drift in rivers, conservative positioning in coastal zones, and energy-conservative strategies in open ocean.

Expected outcomes include 50\% improvement in area mapped per Joule over energy-unaware frontier exploration, 30\% improvement over rule-based energy-aware baselines, mission success rates exceeding 95\% across all water-body types, and guaranteed safe return with at least 10\% battery safety margin. The approach represents the first RL framework for joint ASV motion control and multi-sensor duty-cycling with water-body type adaptation and enforced safe return guarantees.
